% =====================================================================
% 第 VI 章  Remarks on class-invariants and related topics
%          (S. Chowla, Dec. 3 and Dec. 10, 1957)
% =====================================================================
\lecture{Remarks on class-invariants and related topics}{Remarks on Class-Invariants and Related Topics}{S. Chowla}{Dec. 3 and Dec. 10, 1957}

% ---------------------------------------------------------------
\subsection*{\texorpdfstring{§1}{§1}}

Let $p>3$ be a prime. Suppose that the class-number of the imaginary
quadratic field $R(\sqrt{-p})$ is $1$. In this case we have seen that
\[
j\!\left(\frac{-3+\sqrt{-p}}{2}\right)
\tag{1}\label{eq:VI-1}
\]
is a rational integer. Here
\[
j(\omega)=\frac{\bigl\{1+240\sum_{1}^{\infty}\sigma_{3}(n)q^{2n}\bigr\}^{3}}
{q^{2}\prod_{1}^{\infty}(1-q^{2n})^{24}}
\tag{2}\label{eq:VI-2}
\]
and $q=e^{\pi i\omega}$. It is a remarkable fact (see Weber,
\emph{Lehrbuch der Algebra}, Bd.~3, 457--462 \cite{weber1908algebra})
that not only the number \eqref{eq:VI-1} but also its cube-root is a
rational integer. Under the same restriction on the class-number it is
also true (Weber, p.~504) that
\[
\frac{\sqrt{j(\omega)-1728}}{\sqrt{-p}}
\qquad\Bigl\{\omega=\frac{-1+\sqrt{-p}}{2}\Bigr\}
\]
is a rational integer.

We can now formulate the

\begin{namedstmt}{Theorem}
If the class-number of $R(\sqrt{-p})$ is $1$, where $p$ is a prime $>3$,
then the diophantine equation
\[
x^{3}-py^{2}=-1728
\]
has a solution in rational integers with $x=\bigl\{e^{\pi\sqrt{p}/3}
\bigr\}$ where $\{u\}$ denotes the integer nearest to $u$.
\end{namedstmt}

It is believed that $p=7,11,19,43,67,163$ are the only primes $>3$ such
that the class-number of $R(\sqrt{-p})$ is $1$. Heilbronn and Linfoot
have proved that there is at most one more prime $p$ with the property in
question (by D.~H. Lehmer we must have $p>10^{9}$ for such a prime).
Chowla and Selberg (Proc.\ Nat.\ Acad.\ Sci.\ U.S.A.\ 1949)
\cite{chowlaselberg1949} have proved that the existence of a prime
$p>163$ with class-number of $R(\sqrt{-p})$ equal to $1$, would imply the
falseness of the ``extended Riemann hypothesis''. The formulation of the
above theorem is not without interest for the following reason. Siegel's
method (Acta Arithmetica, Bd.~1, 1936) does not lead to a determinable
constant $c$ beyond which the class-number of $R(\sqrt{-p})$ is greater
than $1$. Also the Thue--Siegel--Roth theorem on the finiteness of the
number of solutions of certain diophantine equations (including the form
$ay^{2}=bx^{3}+f$ of our theorem) does not lead to determinable constants
$c$ such that the equations have no solution when the variables in the
diophantine equation are absolutely greater than $c$. Our theorem would
be useful if, for example, one could use the methods of Delaunay and
Nagell to prove the existence of determinable constants
$c_{1}(<\frac{\pi}{3})$ and $c_{2}$ such that $x^{3}-py^{2}=-1728$ has no
solutions in integers for $x>e^{c_{1}\sqrt{p}}$ and $p>c_{2}$. One could
then deduce from our theorem that there is a determinable constant $c_{3}$
such that the class-number of $R(\sqrt{-p})$ is greater than $1$ for
$p>c_{3}$. It should also be remarked here that a recent claim by Kurt
Heegner, (M.Z.\ 1954 \cite{heegner1954}) to have solved this problem,
seems unjustified.

% ---------------------------------------------------------------
\subsection*{\texorpdfstring{§2}{§2}}

If $m$ is a positive integer such that the classes of reduced binary
quadratic forms of discriminant $-m$ have a single class in each genus,
one uses Kronecker's ``Grenz Formel'' to determine explicitly (!) the
values of $j(\omega)$ when $\omega=i\sqrt{m}$.

We introduce the functions
\[
\begin{gathered}
\eta(\omega)=q^{\frac{1}{12}}\prod_{1}^{\infty}(1-q^{2n})\\[3pt]
f(\omega)=e^{\frac{-\pi i}{24}}\,\frac{\eta\bigl(\frac{\omega+1}{2}\bigr)}{\eta(\omega)}
\end{gathered}
\]
(see Weber, p.~113)

Kronecker's formula (``Grenz Formel'') states that (ibid., p.~531)
\[
\begin{aligned}
\lim_{S\to 1+0}&\Bigl\{\Sigma'\bigl(Ax^{2}+2Bxy+Cy^{2}\bigr)^{-S}
-\frac{\pi}{(S-1)\sqrt{m}}\Bigr\}\\
&=-\frac{2\pi\,\Gamma'(1)}{\sqrt{m}}+\frac{\pi}{\sqrt{m}}\log\frac{A}{4m}\\
&-\frac{2\pi}{\sqrt{m}}\log\;\eta(\omega_{1})\,\eta(\omega_{2})
\end{aligned}
\]
where $m=AC-B^{2}$ and
\[
\omega_{1}=\frac{B+i\sqrt{m}}{A},\qquad
\omega_{2}=\frac{-B+i\sqrt{m}}{A}.
\]
(The prime in $\Sigma'$ means that $x=y=0$ is omitted from the summation
where $x,y$ range over all integers.) From this, without difficulty, one
deduces, for example, that
\[
\left(\frac{f(\sqrt{-105})}{\sqrt[4]{2}}\right)^{12}
=(2+\sqrt{5})^{2}\,(55+12\sqrt{21})\,(6+\sqrt{35})\,(2+\sqrt{3})^{3}.
\]

Using the formula
\[
j(\omega)=\frac{\bigl\{f^{24}(\omega)-16\bigr\}^{3}}{f^{24}(\omega)}
\]
we can now deduce the value in algebraic form of $j(\sqrt{-105})$.

% ---------------------------------------------------------------
\subsection*{\texorpdfstring{§3}{§3}}

We add a few remarks on Ramanujan's function $\mathcal{T}(n)$ defined by
\[
\sum_{1}^{\infty}\mathcal{T}(n)\,x^{n}
=x\prod_{1}^{\infty}(1-x^{n})^{24}
\tag{4}\label{eq:VI-4}
\]
($|x|<1$) [cf.\ the denominator of \eqref{eq:VI-2}]. Ramanujan
conjectured and Mordell proved that the function $\mathcal{T}(n)$ is
multiplicative, i.e.
\[
\mathcal{T}(mn)=\mathcal{T}(m)\,\mathcal{T}(n)\quad\text{if }(m,n)=1
\tag{5}\label{eq:VI-5}
\]

Ramanujan also conjectured that for prime $p$,
\[
|\tau(p)|<2p^{\frac{11}{2}}
\tag{6}\label{eq:VI-6}
\]
the best that is known in this direction is that
$\mathcal{T}(p)=O(p^{23/4})$ which follows from A.~Weil's estimate on
Kloosterman's sums
\[
\Bigl|\sum_{x=1}^{p-1}e^{-\frac{2\pi i}{p}(x+\bar{x})}\Bigr|<2\sqrt{p}
\tag{7}\label{eq:VI-7}
\]
Here $\bar{x}$ is defined by $x\bar{x}\equiv 1\pmod{p}$. Another unsolved
problem (raised by D.~H. Lehmer) is whether $\mathcal{T}(n)$ can ever
be $0$.

% ---------------------------------------------------------------
\subsection*{\texorpdfstring{§4}{§4}}

The calculation of class-invariants $j(\omega)$ can also be made to
depend on the classical modular equations of Legendre and Jacobi. We
recall the main details of the connection. For $0<k<1$ write
\[
K=\int_{0}^{\frac{\pi}{2}}\frac{d\phi}{\sqrt{1-k^{2}\sin^{2}\phi}}
=\frac{\pi}{2}\,F\bigl(\tfrac{1}{2},\tfrac{1}{2},1;k^{2}\bigr)
\tag{8}\label{eq:VI-8}
\]
$k'$ is defined by $k^{2}+k'^{2}=1$ and
\[
K'=\int_{0}^{\frac{\pi}{2}}\frac{d\phi}{\sqrt{1-k'^{2}\sin^{2}\phi}}
\tag{9}\label{eq:VI-9}
\]
Also write
\[
q=\exp(-\pi K'/K).
\]
Then
\[
\sqrt{\frac{2K}{\pi}}=\sum_{-\infty}^{\infty}q^{n^{2}}.
\tag{10}\label{eq:VI-10}
\]

Let
\[
L=\int_{0}^{\frac{\pi}{2}}\frac{d\phi}{\sqrt{1-l^{2}\sin^{2}\phi}},\qquad
L'=\int_{0}^{\frac{\pi}{2}}\frac{d\phi}{\sqrt{1-l'^{2}\sin^{2}\phi}}.
\]
If
\[
\frac{L'}{L}=\frac{nK'}{K}
\]
where $n$ is a positive integer, there is an algebraic relation between
$k$ and $l$ (called the ``modular equation''):
\[
n=3:\quad\sqrt{kl}+\sqrt{k'l'}=1
\]
\[
n=7:\quad\sqrt[4]{kl}+\sqrt[4]{k'l'}=1.
\]
Both results are due to Legendre. By setting $k=l'$ (and so $k'=l$) we get
an algebraic equation for $k$ when $\frac{K'}{K}=\sqrt{n}$ and $n$ is a
positive integer. Abel remarked that the resulting equation for $k$ is
``solvable by radicals''.

We quote the following results:
\[
\begin{aligned}
&\text{If }\frac{K'}{K}=\sqrt{210},\text{ then }
k=(\sqrt{2}-1)^{4}(2-\sqrt{3})^{2}(\sqrt{7}-\sqrt{6})^{4}\\
&\qquad\times(8-3\sqrt{7})^{2}(\sqrt{10}-3)^{4}(4-\sqrt{15})^{2}
(\sqrt{15}-\sqrt{14})^{2}\\
&\qquad\times(6-\sqrt{35})^{2}
\end{aligned}
\tag{11}\label{eq:VI-11}
\]
(Ramanujan in a letter to Hardy dated 27 Feb.\ 1913)

\noindent
If $\frac{K'}{K}=\sqrt{47}$, then $\sqrt[12]{16\,kk'}$ is the real root of
\[
x^{5}+2x^{4}+2x^{3}+x^{2}-1=0
\tag{13}\label{eq:VI-13}
\]
(after multiplication by $x-1$ the left-side becomes
$x^{6}+x^{5}-x^{3}-x^{2}-x+1$).

The following solution originated from G.~P.~Young (Am.\ J.\ of Math.\
10, 1888, 99--130) and was simplified by Cayley (ibid., 13, 1891,
53--58). The real root of \eqref{eq:VI-13} is
\[
\sqrt{\frac{\sqrt[5]{A}+\sqrt[5]{B}+\sqrt[5]{C}+\sqrt[5]{D}}{10}}
\]
where
\[
\begin{aligned}
A,\ D&=39000+18200\sqrt{5}\pm(1720+920\sqrt{5})\sqrt{235+94\sqrt{5}}\\
B,\ C&=39000-18200\sqrt{5}\mp(1720-920\sqrt{5})\sqrt{235-94\sqrt{5}}.
\end{aligned}
\]
(See also Journal of the Indian Math.\ Soc.\ 18, 1929--30, p.~273 of
``Notes and questions''.)

% ---------------------------------------------------------------
\subsection*{\texorpdfstring{§5}{§5}}

Finally we quote the following results of Deuring [``Die Typen der
Multiplikatorenringe elliptischer Funktionenk\"orper''] (G.~Herglotz zum
60.~Geburtstag gewidmet) in Abh.\ aus der Math.\ Sem.\ der Hansischen
Univ.\ 14, 1941, 197--272 \cite{deuring1941typen}.

Let
\[
j(\omega)=\frac{\bigl(f^{24}(\omega)-16\bigr)^{3}}{f^{24}(\omega)}
\]
where $\omega=\omega_{1}/\omega_{2}$ is the quotient of two numbers
$\omega_{1}$ and $\omega_{2}$ which form a basis of any ideal in an
imaginary quadratic field $\Sigma$, then by Weber (ibid, 540--541)
\[
f(\omega)\big/\sqrt{2}
\]
is a unit of an algebraic number field if $2$ is ``voll zerlegt'' in
$\Sigma$, otherwise the product of $f(\omega)/\sqrt{2}$ by a certain power
of $2$ with positive exponent is a unit. From this result of Weber (based
on a simple argument using Kronecker's Grenz Formel), Deuring deduces:

A singular invariant $j$ of characteristic $0$ is either not divisible by
any prime factor of $2$ or by every such factor according as $2$ is
``voll zerlegt'' or not in $\Sigma$.

By more difficult arguments Deuring proves analogous results when $2$ is
replaced by $3$, $5$, $7$ or $13$. We quote from p.~271 of his paper
(also p.~257 for a relevant table):

\begin{quote}
$j$ ist durch keinen Primfaktor von $3$ oder $5$ teilbar oder durch
jeden, je nachdem $3$ $(5)$ in $\Sigma$ voll zerlegt oder nicht.

$j-2^{6}\cdot 3^{3}=j+1$ ist durch keinen Primfaktor von $7$ teilbar oder
durch jeden, je nachdem $7$ in $\Sigma$ voll zerlegt oder nicht.

$j-5$ ist durch keinen Primfaktor von $13$ teilbar oder durch jeden, je
nachdem $13$ in $\Sigma$ voll zerf\"allt oder nicht.
\end{quote}

% ---------------------------------------------------------------
\subsection*{\texorpdfstring{§6}{§6}}

In this section we shall give a new proof of the fundamental fact that
the function
\[
j(\tau)=\frac{\bigl\{1+240\sum_{1}^{\infty}\sigma_{3}(n)e^{2n\pi i\tau}\bigr\}^{3}}
{e^{2\pi i\tau}\prod_{1}^{\infty}(1-e^{2n\pi i\tau})^{24}}
\]
is invariant for the transformations $\tau\longrightarrow\tau+1$ and
$\tau\longrightarrow -\frac{1}{\tau}$ (the first is, of course, obvious).
Here $\operatorname{Im}(\tau)>0$ and
\[
\sigma_{a}(n)=\sum_{d\mid n}d^{a}.
\]

The proof is based on the functional equation for the Riemann zeta
function
\[
2\,\Gamma(s)\,\zeta(s)\cos\frac{s\pi}{2}=(2\pi)^{s}\zeta(1-s)
\tag{1}\label{eq:VI2-1}
\]
and on Mellin's Integral ($s=\sigma+it$)
\[
\frac{1}{2\pi i}\int_{(\sigma=k>0)}\Gamma(s)y^{-s}\,ds=e^{-y}.
\tag{2}\label{eq:VI2-2}
\]
Here $\operatorname{Re}(y)>0$ and the path of integration is the vertical
line $\sigma=k$ where $k>0$. We first prove the

\begin{namedstmt}{Theorem}
\emph{Write
\[
H_{a}(y)=\sum_{1}^{\infty}\sigma_{a}(n)e^{-2n\pi y}
+(-1)^{\frac{a+1}{2}}\frac{\Gamma(a+1)\zeta(a+1)}{(2\pi)^{a+1}}.
\tag{3}\label{eq:VI2-3}
\]}
Then if $a$ is an odd integer $\geq 3$, we have
\[
(-1)^{\frac{a+1}{2}}y^{a+1}H_{a}(y)=H_{a}\bigl(1/y\bigr).
\tag{4}\label{eq:VI2-4}
\]
\end{namedstmt}

There is a slightly different formulation for the case $a=1$, which
enables us to study the effect of the transformation
$\tau\longrightarrow -\frac{1}{\tau}$ on the denominator of $j(\tau)$.
For the numerator we use our theorem for the case $a=3$. Combining these
two results we deduce the property of $j(\tau)$ in question. Write
\[
G(y)=G_{a}(y)=\sum_{1}^{\infty}\sigma_{a}(n)e^{-2n\pi y}
\tag{5}\label{eq:VI2-5}
\]
and
\[
F(s)=\zeta(s)\,\zeta(s-a)\quad\Bigl[=\sum_{1}^{\infty}\frac{\sigma_{a}(n)}{n^{s}}
\text{\ \ for\ \ }\sigma>a+1\Bigr].
\tag{6}\label{eq:VI2-6}
\]

From \eqref{eq:VI2-2} and \eqref{eq:VI2-6}
\[
G(y)=\frac{1}{2\pi i}\int_{(\sigma=a+2)}\Gamma(s)F(s)(2\pi y)^{-s}\,ds.
\tag{7}\label{eq:VI2-7}
\]

From \eqref{eq:VI2-1}
\[
2\,\Gamma(s-a)\,\zeta(s-a)\cos\frac{(s-a)\pi}{2}=(2\pi)^{s-a}\zeta(1-s+a).
\tag{8}\label{eq:VI2-8}
\]

Since $a$ is odd it follows from \eqref{eq:VI2-6}, \eqref{eq:VI2-1} and
\eqref{eq:VI2-8} that
\[
\begin{aligned}
2\,\Gamma(s)&\Gamma(s-a)\,F(s)\,\sin(s\pi)\,(-1)^{\frac{a-1}{2}}\\
&=(2\pi)^{2s-a}F(a+1-s).
\end{aligned}
\tag{9}\label{eq:VI2-9}
\]

Multiplying both sides by $\Gamma(a+1-s)$ and using
$\Gamma(x)\Gamma(1-x)=\pi/\sin(\pi x)$ we get
\[
\Gamma(s)F(s)(-1)^{\frac{a+1}{2}}
=(2\pi)^{2s-a-1}\,\Gamma(a+1-s)\,F(a+1-s).
\tag{10}\label{eq:VI2-10}
\]

From \eqref{eq:VI2-7} and \eqref{eq:VI2-10}
\[
\begin{aligned}
G(y)&=\frac{1}{2\pi i}\int_{(\sigma=a+2)}
(2\pi)^{2s-a-1}(2\pi y)^{-s}(-1)^{\frac{a+1}{2}}
\Gamma(a+1-s)F(a+1-s)\,ds\\
&=\frac{1}{2\pi i}\int_{t=-\infty}^{t=+\infty}
(2\pi)^{a+3+2ti}(2\pi y)^{-a-2-ti}(-1)^{\frac{a+1}{2}}
\Gamma(-1-ti)F(-1-ti)\,dt
\end{aligned}
\]
(and now transform by $t\to -t$)
\[
\begin{aligned}
&=\frac{1}{2\pi i}\int_{t=-\infty}^{t=+\infty}
(2\pi)^{a+3-2ti}(2\pi y)^{-a-2+ti}(-1)^{\frac{a+1}{2}}
\Gamma(-1+ti)F(-1+ti)\,dt\\
&=\frac{1}{2\pi i}\int_{(\sigma=-1)}
(2\pi)^{a+1-2s}(2\pi y)^{-a-1+s}(-1)^{\frac{a+1}{2}}
\Gamma(s)F(s)\,ds
\end{aligned}
\]
so that
\[
G(y)=(-1)^{\frac{a+1}{2}}\left(\frac{1}{y}\right)^{a+1}
\int_{(\sigma=-1)}\left(\frac{y}{2\pi}\right)^{s}\Gamma(s)F(s)\,ds.
\tag{11}\label{eq:VI2-11}
\]

Next, from \eqref{eq:VI2-7},
\[
\begin{aligned}
G(1/y)&=\frac{1}{2\pi i}\int_{(\sigma=a+2)}\left(\frac{y}{2\pi}\right)^{s}
\Gamma(s)F(s)\,ds\\
&=\frac{1}{2\pi i}\int_{(\sigma=-1)}\left(\frac{y}{2\pi}\right)^{s}
\Gamma(s)F(s)\,ds\\
&\quad+\Bigl\{\text{sum of residues of integrand at poles }
s=a+1,\ s=1,\ s=0\Bigr\}.
\end{aligned}
\]
Thus, using \eqref{eq:VI2-11},
\[
\begin{aligned}
G(1/y)&=(-1)^{\frac{a+1}{2}}\,y^{a+1}\,G(y)\\
&+\Bigl\{\Gamma(a+1)\,\zeta(a+1)\left(\frac{y}{2\pi}\right)^{a+1}\\
&\qquad+\Gamma(1)\,\zeta(1-a)\left(\frac{y}{2\pi}\right)\\
&\qquad+\zeta(0)\,\zeta(-a)\Bigr\}.
\end{aligned}
\tag{12}\label{eq:VI2-12}
\]

Since $a$ is odd and $\geq 3$ it follows from \eqref{eq:VI2-1} that
$\zeta(1-a)=0$, and so we have from the last relation:
\[
\begin{aligned}
G(1/y)&=(-1)^{\frac{a+1}{2}}\,y^{a+1}\,G(y)\\
&\quad+\Gamma(a+1)\,\zeta(a+1)\left(\frac{y}{2\pi}\right)^{a+1}
+\zeta(0)\,\zeta(-a).
\end{aligned}
\]

Using \eqref{eq:VI2-1} and $\zeta(0)=-\frac{1}{2}$ this gives for odd
$a\geq 3$,
\[
\begin{aligned}
(-1)^{\frac{a+1}{2}}\,y^{a+1}&\Bigl\{G_{a}(y)
+(-1)^{\frac{a+1}{2}}\frac{\Gamma(a+1)\zeta(a+1)}{(2\pi)^{a+1}}\Bigr\}\\
&=G_{a}(1/y)+(-1)^{\frac{a+1}{2}}\frac{\Gamma(a+1)\zeta(a+1)}{(2\pi)^{a+1}}
\end{aligned}
\]
and hence, from \eqref{eq:VI2-3}, we obtain \eqref{eq:VI2-4} which is our
theorem. We have to modify our theorem slightly for the case $a=1$. In
fact when $a=1$, \eqref{eq:VI2-12} gives
\[
\begin{aligned}
G_{1}(1/y)&=-y^{2}G_{1}(y)+\frac{1}{24}y^{2}-\frac{y}{4\pi}\\
&\quad+\frac{1}{2}\cdot\frac{2\,\Gamma(2)\,\zeta(2)}{(2\pi)^{2}},
\end{aligned}
\]
or
\[
-\frac{1}{24}+G_{1}(1/y)
=-y^{2}\Bigl\{G_{1}(y)-\frac{1}{24}\Bigr\}-\frac{y}{4\pi}.
\tag{13}\label{eq:VI2-13}
\]

From \eqref{eq:VI2-13} we shall deduce that
\[
\frac{e^{-2\pi y}\prod_{1}^{\infty}(1-e^{-2n\pi y})^{24}}
{e^{-2\pi/y}\prod_{1}^{\infty}(1-e^{-2n\pi/y})^{24}}=y^{-12}
\tag{14}\label{eq:VI2-14}
\]
which shows the effect of the transformation
$\tau\longrightarrow -\frac{1}{\tau}$ on the denominator of $j(\tau)$. To
prove \eqref{eq:VI2-14} we write
\[
f(q)=q\prod_{1}^{\infty}(1-q^{n})^{24}=g(\tau)
\tag{15}\label{eq:VI2-15}
\]
where
\[
q=e^{2\pi i\tau}.
\tag{16}\label{eq:VI2-16}
\]

We have by ``logarithmic differentiation'',
\[
q\,\frac{f'(q)}{f(q)}=1-24\sum_{1}^{\infty}\sigma_{1}(n)e^{2n\pi i\tau}
\tag{17}\label{eq:VI2-17}
\]
\[
f'(q)=g'(\tau)\,\frac{d\tau}{dq}=\frac{g'(\tau)}{2\pi iq}.
\]
So
\[
q\,\frac{f'(q)}{f(q)}=\frac{g'(\tau)}{2\pi i\,g(\tau)}.
\tag{18}\label{eq:VI2-18}
\]

From \eqref{eq:VI2-13},
\[
\begin{aligned}
\tau^{2}&\Bigl\{1-24\sum_{1}^{\infty}\sigma_{1}(n)e^{2n\pi i\tau}\Bigr\}\\
&=\Bigl\{1-24\sum_{1}^{\infty}\sigma_{1}(n)e^{-\frac{2n\pi i}{\tau}}\Bigr\}
+\frac{6i\tau}{\pi},
\end{aligned}
\tag{19}\label{eq:VI2-19}
\]
or, using \eqref{eq:VI2-17}, \eqref{eq:VI2-18}, \eqref{eq:VI2-19},
\[
\tau^{2}\,\frac{g'(\tau)}{g(\tau)}\cdot\frac{1}{2\pi i}
=\frac{g'(-\frac{1}{\tau})}{2\pi i\,g(\frac{1}{\tau})}
+\frac{6i\tau}{\pi}.
\tag{20}\label{eq:VI2-20}
\]

Hence
\[
\frac{g'(\tau)}{g(\tau)}
=\frac{1}{\tau^{2}}\,\frac{g'(-\frac{1}{\tau})}{g(-\frac{1}{\tau})}
-\frac{12}{\tau}
\tag{21}\label{eq:VI2-21}
\]
integrating,
\[
\log g(\tau)=\log g\!\left(-\frac{1}{\tau}\right)-12\log\tau+\log k
\tag{22}\label{eq:VI2-22}
\]
where $k$ is a constant independent of $\tau$. Thus
\[
\frac{g(\tau)}{g(-\frac{1}{\tau})}=\frac{k}{\tau^{12}}.
\tag{23}\label{eq:VI2-23}
\]

Setting $\tau=iy$ this becomes
\[
\frac{e^{-2\pi y}\prod_{1}^{\infty}(1-e^{-2n\pi y})^{24}}
{e^{-2\pi/y}\prod_{1}^{\infty}(1-e^{-2n\pi/y})^{24}}=\frac{k}{y^{12}}.
\]
Set $y=1$ here to get $k=1$; \eqref{eq:VI2-14} is proved. From
\eqref{eq:VI2-23} [or \eqref{eq:VI2-14}] and the case $a=3$ of our theorem
[\eqref{eq:VI2-3} and \eqref{eq:VI2-4}] we see at once that
$j(\tau)=j(-1/\tau)$.

% ---------------------------------------------------------------
\subsection*{\texorpdfstring{§7}{§7}}

By combining the results of Deuring in §5 with our theorem in §1 we get
the:

\begin{namedstmt}{Theorem}
If $p$ is a prime $>19$ and if the class-number of $R(\sqrt{-p})$ is $1$,
the equation
\[
x^{3}-py^{2}=-8
\]
has a solution in rational integers and indeed with
$x=\frac{1}{6}\bigl\{e^{\frac{\pi}{3}\sqrt{p}}\bigr\}$ where $\{u\}$
denotes the integer nearest to $u$.
\end{namedstmt}

Incidentally we observe that if $p>19$ is a prime, and the class-number
of $R(\sqrt{-p})$ is $1$, then
\[
\Bigl\{e^{\frac{\pi}{3}\sqrt{p}}\Bigr\}
\]
is a multiple of $6$!

% ---------------------------------------------------------------
\subsection*{\texorpdfstring{§8}{§8}}

Hecke [\"Uber die Kroneckersche Grenzformel f\"ur reelle quadratische
K\"orper und die Klassenzahl relativ Abelscher Zahlk\"orper.
Verhandlungen der Naturforschenden Gesellschaft in Basel, Bd.\ 28$_{2}$,
1917, S.~363 \cite{hecke1917kroneckersche}; see also his paper
``Bestimmung der Klassenzahl einer neuen Reihe von algebraischen
Zahlk\"orpern'', G\"ott.\ Nachr.\ 1921, S.~1 \cite{hecke1921bestimmung}]
has proved that for a real quadratic field of discriminant $\Delta>0$,
the Grenzformel takes the following form: let $\xi>1$ be a unit of the
field:
\[
\varepsilon=\frac{u+v\sqrt{\Delta}}{2}\qquad(u,v\geq 1)
\]
and if $\alpha$ runs through a complete system of non-associated (with
respect to $\pm\varepsilon^{v}$) numbers of the ideal
\[
\mathfrak{a}=[\alpha_{1},\alpha_{2}],\qquad
\alpha_{1}\alpha_{2}'-\alpha_{2}\alpha_{1}'=N\sqrt{\Delta}>0,\quad
N=N(\mathfrak{a}),
\]
then we have
\[
\begin{aligned}
N(\mathfrak{a})^{s}\sum_{\alpha}\frac{1}{|N(\alpha)|^{s}}
&=\frac{\lg\varepsilon^{2}}{\sqrt{\Delta}}\cdot\frac{1}{s-1}\\
&\quad+\frac{2}{\sqrt{\Delta}}\,\phi(v)+\theta_{1}(s-1)
+\theta_{2}(s-1)^{2}+\dots
\end{aligned}
\]
where
\[
\phi(v)=E\lg\varepsilon^{2}
-\int_{-\lg\varepsilon}^{+\lg\varepsilon}
\lg\!\left(\sqrt[4]{\Delta\sqrt{\frac{\omega\bar{\omega}}{2i}}}\;
\eta(\omega)\,\eta(-\bar{\omega})\right)dv
\]
where $E$ is Euler's constant and
\[
\omega=-\frac{\alpha_{2}e^{v}-i\alpha_{2}'}{\alpha_{1}e^{v}-i\alpha_{1}'},
\qquad
\bar{\omega}=-\frac{\alpha_{2}e^{v}+i\alpha_{2}'}{\alpha_{1}e^{v}+i\alpha_{1}'}.
\]

Herglotz (Ber.\"u.\ a.\ Verhandlungen.\ Akad.\ Wiss.\ Leipzig.\ Math.\
Phys.\ Klasse, 75, 1923, 3--14 \cite{herglotz1923bestimmung}) obtains
some curious results from Hecke's formula. The following is a sample:
\[
\begin{aligned}
\varepsilon=4+\sqrt{15},\qquad
\int_{0}^{1}\frac{\lg(1+t^{\varepsilon})}{1+t}\,dt
&=-\frac{\pi^{2}}{12}\bigl(\sqrt{15}-2\bigr)\\
&\quad+\lg 2\,\lg(\sqrt{3}+\sqrt{5})
+\lg\frac{1+\sqrt{5}}{2}\,\lg(2+\sqrt{3}).
\end{aligned}
\]

A direct evaluation of this definite integral is probably difficult!
